% SPDX-License-Identifier: MIT-0
\documentclass[11pt]{article}

\usepackage[T1]{fontenc}
\usepackage[utf8]{inputenc}
\usepackage{lmodern}
\usepackage[a4paper,margin=2.4cm]{geometry}
\usepackage{amsmath,amssymb}
\usepackage{microtype}
\usepackage{booktabs}
\usepackage{longtable}
\usepackage[dvipsnames]{xcolor}
\usepackage{titlesec}
\usepackage[colorlinks=true,linkcolor=MidnightBlue,urlcolor=MidnightBlue,
            citecolor=MidnightBlue]{hyperref}

\hypersetup{%
  pdftitle={The surreal and surcomplex reports},
  pdfauthor={Vladimir Reshetnikov},
  pdfsubject={Catalogue of eighteen research packages on surreal and surcomplex numbers}}

\definecolor{Ink}{HTML}{16324F}
\definecolor{Accent}{HTML}{4E5C3D}

\titleformat{\section}{\Large\bfseries\color{Ink}}{\thesection}{0.7em}{}
\titleformat{\subsection}{\large\bfseries\color{Accent}}{\thesubsection}{0.7em}{}

\newcounter{entrynum}
\setcounter{entrynum}{0}

\emergencystretch=1.6em

\def\UrlBreaks{\do\/\do\-\do\_\do\.}
\def\UrlBigBreaks{\do\/}

% \entry{title}{directory}{main document}{source archive}{description}
\newcommand{\entry}[5]{%
  \par\medskip\noindent
  \refstepcounter{entrynum}%
  \textbf{\theentrynum.\; #1}\par\nopagebreak
  \nopagebreak{\raggedright\small\color{Ink}%
    \nolinkurl{#2/}\textbf{\nolinkurl{#3}}\par}\nopagebreak
  \nopagebreak{\raggedright\footnotesize\color{gray}%
    from \texttt{#4}\par}\nopagebreak
  \smallskip\noindent #5\par}

\newcommand{\dir}[1]{\texttt{#1}}
\newcommand{\Q}{\mathbb{Q}}
\newcommand{\R}{\mathbb{R}}
\newcommand{\N}{\mathbb{N}}
\newcommand{\C}{\mathbb{C}}
\newcommand{\No}{\mathbf{No}}
\newcommand{\SC}{\mathbf{No}[i]}
\newcommand{\Z}{\mathbb{Z}}
\newcommand{\der}{\partial}


\begin{document}

\begin{center}
 {\LARGE\bfseries\color{Ink} The surreal and surcomplex reports}\\[8pt]
 {\large Eighteen independent mathematical research packages,\\
 on the surreal field, its algebraic closure,\\
 and the foundations and computation of both}\\[10pt]
 {\small September 2026}
\end{center}

\vspace{4pt}

\section*{Scope}
\addcontentsline{toc}{section}{Scope}

This document catalogues the eighteen research packages stored under
\dir{docs}.  Each entry names the problem the report attacks, the result it
claims, and the archive it arrived in.

The collection divides in three, and so does the directory tree.  Six reports,
under \dir{docs/surreal}, are about the surreal field $\No$ itself: birthdays
and their behaviour under multiplication, the option graphs of surreal forms,
genetic functions and gaps, and two questions of Lipparini about evaluating
formal sums.  Ten, under \dir{docs/surcomplex}, are about the algebraic
closure $\SC$: the theory of holomorphic functions over it, its differential
algebra, its rank-one analytic geometry, and its finite matrices.  The last two,
under \dir{docs/foundations-and-computation}, stand outside both: they are
about the subject rather than inside it, asking in what foundation this
mathematics legally lives and how one computes with it exactly.

These reports were previously part of a larger collection, in a repository
whose subject was elsewhere; they were moved here with their git history.
That is why an entry may refer to a merge, a duplicate sweep or a delivery
number: those describe how the report reached its present form, and the
record is kept rather than rewritten.

Each directory holds a typeset article with its \LaTeX{} source, a
\texttt{README.md}, and in most cases verification code together with the
recorded output of running it.  A merged report additionally keeps every
source manuscript verbatim under \texttt{sources/}, so the merge can be
audited against what went into it.  These are AI-assisted drafts.  None is
refereed or machine-checked, and this catalogue records what each report
claims rather than verifying it.

\medskip
\noindent
Rebuild this document with
\begin{quote}\ttfamily
latexmk -pdf manifest.tex \&\& latexmk -c manifest.tex
\end{quote}

\section{The surreal field}

Six reports about $\No$.  Two questions of Lipparini on evaluating formal
sums; two independent halves of Gonshor's product-birthday conjecture, on
domains neither of which contains the other; the subgraph question from
Roughan's paper on birthday arithmetic; and a genetic counterexample
separating pointwise derivatives from differences of constants.  The first
and the fourth each arrived as a pair of archives and were merged.

\entry{A square-root obstruction to evaluating increasing Hahn series at $\omega$}
{surreal/hahn-evaluation-at-omega}
{article.pdf}{hahn\_evaluation\_research.zip + reversed\_hahn\_series.zip}
{Answers Problem~7.7 (p.~38) of Lipparini (arXiv:2505.00424v2) negatively: there
is no ring homomorphism from the field of set-supported real-coefficient Hahn
series with strictly increasing surreal exponents into the surreals sending
$x^a$ to $\omega^a$.  The certificate is $B\in\Q[[x]]$ with $B^2=1-x$, built by
the recursion $b_n=-\tfrac12\sum_{k=1}^{n-1}b_kb_{n-k}$: any such map would force
a square to equal $1-\omega<0$.  The homomorphism is applied only to the finite
identity, so no continuity, valuation, order preservation or preservation of
infinite sums is assumed.  Adds the exact classification of admissible images of
the variable as the infinitesimals, the monomial classification over an
arbitrary ordered abelian group, and order-uniqueness on the source field.
\emph{Merged; the core certificate was identical in both originals and is proved
once, but the two toolkits are kept --- in particular a nonnegative-coefficient
semiring obstruction with strictly weaker hypotheses that the square-root
certificate cannot reach.}}

\entry{Gonshor's Product-Birthday Conjecture for ordinal-support series}
{surreal/gonshor-product-birthdays}
{surreal\_product\_birthdays.pdf}{gonshor\_product\_birthdays\_research.zip}
{Proves Gonshor's bound $b(xy)\le b(x)\otimes b(y)$ on the ring of Conway normal
forms $\sum_{\alpha\in S}c_\alpha\omega^{-\alpha}$ with $S$ an arbitrary set of
ordinals, so transfinite and uncountable supports are allowed.  Supplies an
exact birthday formula by transfinite induction over the support with maintained
prefix bounds and an explicit limit stage, and a complete classification of the
equality cases.  The route is a leading-Cantor-degree gap, which also gives that
the degree of the birthday never rises under multiplication.  The conjecture for
arbitrary surreal numbers is not solved.  The domain is incomparable with that of the
Laurent-series report that follows: it excludes $\omega$ and every positive
power.}

\entry{A Laurent-series case of Gonshor's birthday conjecture}
{surreal/gonshor-laurent-birthdays}
{article.pdf}{Gonshor\_Laurent\_Birthday\_Research.zip}
{Proves the same bound $b(xy)\le b(x)\otimes b(y)$ on the canonically embedded
field $\R((\omega^{-1}))$, which unlike the ring of the preceding entry
contains $\omega$
and its positive powers but restricts supports to order type at most $\omega$ on
the integer lattice.  The route is a bilinear budget: an exact splitting
$b(x)=b(P)\oplus b(u)$ into polynomial and infinitesimal parts, three separate
product estimates for the resulting cross terms, and a finite-jet stabilisation
carrying the finite bound to infinite series without any continuity claim.  Also
treats finite normal forms with arbitrary ordinal exponents, derives exact
birthdays for finite and infinite Laurent series, and proves a reciprocal
corollary.  Neither this theorem nor the preceding one contains the other; on
their common corner $\R[[\omega^{-1}]]$ the two agree and are proved twice.}

\entry{Canonical forms need not be subgraphs}
{surreal/canonical-forms-need-not-be-subgraphs}
{surreal\_graphs.pdf}{surreal\_graphs\_research.zip + surreal\_sparse\_graphs.zip}
{Answers negatively the question in the Figure~3 caption of Roughan,
\emph{Surreal Birthdays and Their Arithmetic} (arXiv:1810.10373v2, p.~7):
the canonical form of a surreal value need not be a subgraph of every other form
of that value.  Five vertices and rank three are both least possible, and there
are exactly sixteen five-vertex counterexamples, eight triangle-free, four of
value $1/2$.  Adds cycle representations of $1/2$ at every admissible length,
the exact extremal function, a planar construction of maximum degree three and
arbitrarily large girth, classifications of which canonical graphs are
unavoidable, and transfinite witnesses.  \emph{Merged; the two originals
exhibited genuinely different five-vertex witnesses --- one passing through a
vertex of value $-1$, one with no negative value at all --- and both are printed
in full, as are two odd-cycle families, two high-girth constructions and two
proofs of the Fibonacci bound.}}

\entry{A countable genetic cut across the finite--infinite gap}
{surreal/genetic-gaps-and-primitives}
{article.pdf}{genetic\_gap\_counterexamples.zip}
{With $q(t)=t/(1+t^2)$ and $B(x)=1-\{q(n-x):n\in\N\mid\;\}$ as a Conway cut,
proves that $B$ is the indicator of the positive infinite surreals, is genetic
under the set-sized option constructors of the cited sources, and has pointwise
derivative zero everywhere.  Consequently $x$ and $x-B(x)$ are genetic
primitives of $1$ agreeing at every real input that do not differ by a constant,
the latter also an increasing bijection.  The sublevel class
$\{x:0<x<\omega,\;B(x)\le\tfrac12\}$ is nonempty and bounded yet has no surreal
supremum, so neither end gap supplies one.}

\entry{A Number-Valued Broadcast Sum of Surreal Sequences}
{surreal/broadcast-sum-of-surreal-sequences}
{article.pdf}{surreal\_broadcast\_sum.zip}
{Analyses the sign-truncation game of Lipparini (arXiv:2505.00424v2,
Remark~7.6(3), in connection with Problem~7.5).  Proves the game always has a
surreal-number value, that it agrees with finite surreal sums and with
Lipparini's ordinal operation on their respective domains, and that it fails
weak coordinatewise monotonicity, with an explicit family
$s[n]=3\cdot2^{-n-d-1}$ on a subset $E$ and $2^{-n-d}$ off it realising the
failure.  The underlying rule is Lipparini's; the name and the notation are
local to the report.}

\section{The surcomplex plane}

Ten reports about $\SC=\No[i]$.  Seven are about holomorphic function theory,
rest on one framework, and are all merges: twenty-seven manuscripts reduced to
seven documents.  The last three came later and are listed after them --- a
merge of seven manuscripts on differential algebra, and two standalone reports,
on rank-one analytic geometry and on finite matrices, each written against this
repository and each naming the gap it fills.

\medskip
\noindent
\emph{Read \dir{analysis} first.}  It is the foundation the other six
continue, and it is where the shared calculus is established: there is no
topological convergence anywhere in this theory.  Every \emph{set} of
surcomplex numbers is closed and discrete, a convergent set-indexed net is
eventually constant, and every continuous map $[0,1]\to\SC$ is constant, so
$\sum t^n=(1-t)^{-1}$ is a Hahn identity and never a limit of partial sums,
and contour integration is defined coefficientwise on ordinary complex paths.

\medskip
\noindent
\emph{Then read the ring table in \dir{analytic-geometry}, before anything
else.}  Four coefficient rings circulate through this material under nearly
identical notation --- the radius-free $\C\{z\}((t^\Gamma))$, the
fixed-polydisk $\mathcal{O}(D)((t^\Gamma))$, the common-domain germ ring, and
the formal $\C[[z]]((t^\Gamma))$ --- and the inclusions between them are
strict, with $F(z)=\sum_{r\ge1}t^{r\eta}/(1-rz)$ the witness: its coefficient
at $r\eta$ has a pole at $1/r$, so no ordinary disk carries the whole family.
A theorem true over one of these rings can be false over another.  Every
report therefore names the ring of every statement, and records that no
theorem was transported between rings; the one place a bridge is used, it is
one-directional and flagged in the proof that uses it.  The remaining five
reports can be read in any order.

\entry{Surcomplex analysis: holomorphic functions on $\SC$}
{surcomplex/analysis}
{article.pdf}{surcomplex\_analysis.zip and eight companion archives}
{One article assembled from nine manuscripts, all delivered under the same
name on 21 September 2026 as independent attempts at one subject: a theory of
holomorphic functions over $\mathbf{K}=\SC$ built on Hahn-supported
series.  On the convergence-free base described above the article develops
germs, the coherent algebra $\mathcal H(U)$ over an ordinary domain,
Weierstrass preparation and infinitesimal zero clusters, a residue calculus,
the global principles, exponentials and logarithms, and an all-scale rigidity
theorem.  \textbf{Nine independent runs at a foundation do not agree the way
nine runs at a theorem would, and the merge turned on that.}  Of seventeen
load-bearing notions, six are genuinely inequivalent across the sources, and a
theorem proved under one is false under another: there are three inequivalent
function classes (the identity theorem, the global maximum principle,
Liouville and the open mapping theorem hold for coherent sections and fail for
germs; sharp Schwarz and Schwarz--Pick hold for lifts and fail for coherent
sections), and three inequivalent residues, two of which disagree numerically
--- for $1/(\zeta-\varepsilon)$ with $\varepsilon$ infinitesimal the cluster
residue at the ordinary centre $0$ is $1$ while the local residue at $0$ is
$0$, the actual pole sitting at $\varepsilon$.  A single ordinary bound in
Liouville gives a characterization and not constancy: $t z$ is real-bounded by
$1$ on the whole finite plane and is not constant.  Eleven further conflicts
are resolved and recorded, including two functions both called ``the''
logarithm that differ by exactly $2\pi i$.  All nine manuscripts and all nine
verification programs are kept; the deduplication is in the article.  Its
\nolinkurl{MERGE_NOTES.md} is worth reading before the article.}

\entry{Infinitesimal analytic geometry over the surcomplex numbers}
{surcomplex/analytic-geometry}
{article.pdf}{Surcomplex\_Analytic\_Geometry.zip + surcomplex\_finite\_geometry.zip}
{The ring theory of the infinitesimal neighbourhood --- the monad
$\mathfrak{m}_K^n$ of an ordinary point --- rather than the deformation theory
of the reports that follow.  Proves, \emph{separately of each of} the
radius-free $A_n=\C\{z_1,\dots,z_n\}((t^\Gamma))$ and the common-domain
$R_n$, that it is a Noetherian Jacobson domain of dimension $n$ whose maximal
ideals are exactly the infinitesimal evaluation ideals
$\mathfrak{m}_a=(z_1-a_1,\dots,z_n-a_n)$ with residue field $K$, regular local
at each, with completion $K[[X_1,\dots,X_n]]$ --- and that neither is local,
each recording every infinitesimal point at once.  Carries Weierstrass
division and preparation, a weak and a strong Nullstellensatz, Noether
normalization, finite projection, contour residue duality, Hensel--Rouch\'e,
and workspace invariance: every common zero over any larger divisible
workspace already lies in $K_\Gamma^n$ with unchanged multiplicity, so the
count accounts for all zeros in the whole proper class $\SC$.  The
nonzero-value-group hypothesis is essential and not aesthetic --- for
$\Gamma=\{0\}$ the strong Nullstellensatz fails at $I=(0)$.  \emph{This is the
report that fixes the four-ring vocabulary and proves the inclusions strict;
statements holding over both $A_n$ and $R_n$ are stated twice with both
proofs, since the two proofs use different resources and neither transfers.}}

\entry{Finite Hahn deformations: division, conservation of multiplicity, residue duality}
{surcomplex/finite-deformations}
{article.pdf}{surcomplex\_finite\_deformations.zip, \_finite\_fibres, \_finite\_maps,
\_intersections, \_research}
{Five manuscripts on one theorem: an ordinary holomorphic complete
intersection $f$ on a fixed polydisk, with isolated zero at the origin and
colength $\mu$, keeps its entire finite geometry under any Hahn perturbation
$F=f+E$ of well-ordered \emph{strictly positive} error support --- on the
unshrunk domain, with an explicit support certificate
$\mathrm{supp}\subseteq\mathrm{supp}(G)+M$, $M$ the monoid generated by the
error support.  The remainder of division is unique; the witnesses are
explicitly not, their ambiguity being a Koszul syzygy carried by the same
certificate.  Gives the integral structure $H^+_\Gamma(D)/(F)\cong(K^\circ)^\mu$
over a valuation ring that may have high-rank value group and need not be
Noetherian, with nothing completed; conservation of multiplicity at the
\emph{actual} displaced surcomplex zeros rather than their standard parts,
$\sum_\zeta\mu_\zeta=\mu$; and a coefficientwise residue $\mathrm{Res}_F$ ---
defined by a series, not by a surcomplex contour integral --- whose Gram
determinant is a unit even at a collision and which satisfies
$\mathrm{Res}_F(G\cdot J_F)=\mathrm{Tr}(M_G)$, individual summands routinely
infinite while the total is finite.  Also Hartogs, characters and local
algebras, spectral invariants, trace-Jacobian and B\'ezoutian, monodromy and
stability.  The engine throughout is positive-support Neumann inversion; no
argument uses convergence of partial sums in any topology.}

\entry{Jordan separation, Cauchy theory, and Stokes calculus on the surcomplex plane}
{surcomplex/contours-and-stokes}
{article.pdf}{surcomplex\_contours.zip and two same-named archives}
{Three manuscripts whose titles differ by one word --- Stokes \emph{Calculus},
\emph{Integration}, \emph{Formulae} --- over one body of work.  All three
refuse to transfer the Jordan--Cauchy--Stokes package to $\SC$ as a unit, and
all three split it the same way: a coarse Jordan separation theorem in the
standard-part topology whose separator is the whole boundary monad rather than
a curve; an imported semialgebraic Jordan theorem; a Hahn-coefficient de Rham
complex with coefficientwise Stokes, extended to infinitesimally deformed
chains; and the holomorphic consequences --- Cauchy, endpoint displacement,
integer winding, and microscopic circle and torus realizations of the
perturbation residue --- with an algebraic rational contour pairing
characterized uniquely by vanishing on exact differentials and normalizing
$dz/(z-a)$ to $2\pi i\cdot\mathrm{wind}_{\mathrm{sa}}$.  Adds a purely
algebraic third chain category over any characteristic-zero field, a
support-preserving Poincar\'e lemma, an order-valued ML estimate, and a
root-of-unity/Schnirelman section.  This is where the question the
deformation report leaves open is answered: supplying an actual surcomplex
contour that represents the residue series.  \emph{Two thirds of each
manuscript was the same mathematics in the same order; 82 numbered results
became 59, by duplicate elimination only.}}

\entry{Global support obstructions: moving divisors, Mittag--Leffler, and a Picard dichotomy}
{surcomplex/global-divisors}
{article.pdf}{surcomplex\_global\_support.zip + surcomplex\_global\_theorems.zip}
{The local-to-global theory: when do infinitely many individually legitimate
infinitesimal prescriptions admit one coherent global realization over the
ordinary complex plane?  The sharp answer for divisors is a
\emph{symmetric-support} criterion --- a cluster family
$P_c(u)=u^{m_c}+\sum_{j<m_c}a_{c,j}u^j$ over a closed discrete ordinary
$A\subset\C$ is the exact zero divisor of a nonzero global section iff
$\bigcup_c\bigcup_j\mathrm{supp}(a_{c,j})$ is well ordered: the supports of
the symmetric coefficients, not of separately chosen roots.  Hence the two
discriminating examples --- $n+t^{1/n}$ alone is not realizable, its
descending symmetric support $\{1/n\}$ failing, while \emph{all} roots of
$(z-n)^n=t$ are, symmetric support the singleton $\{1\}$ --- and hence a
globally principal effective divisor with a nonprincipal effective
sub-divisor.  Also Mittag--Leffler at fixed centres and for moving poles, a
support-restricted Chinese remainder theory with an explicit Neumann right
inverse, an exact Cousin criterion in terms of singular support, and the
Picard dichotomy: every ordinary-base stalk is a nonlocal principal ideal
domain, $\mathrm{Pic}(\C,\mathcal{H}_\Gamma)=0$ iff $\Gamma=\{0\}$ or
$\Gamma=\mathbb{Z}\delta$, and otherwise $\dim_\C\mathrm{Pic}\ge2^{\aleph_0}$,
with equality for countable noncyclic $\Gamma$.}

\entry{Surcomplex polynomial algebra: root geometry, Newton profiles, residue duality}
{surcomplex/polynomial-algebra}
{article.pdf}{surcomplex\_polynomial\_algebra.zip and two same-named archives}
{Finite-degree polynomial algebra over $\SC$, worked inside set-sized
algebraically closed Hahn workspaces and keeping two geometries deliberately
apart throughout: the surreal-valued modulus, whose disks are ordered, and the
natural Hahn valuation, whose balls have arbitrary ordered-group rank.  A
Rouch\'e theorem in each.  Covers the finite polynomial algebra and the
fundamental theorem there, CRT/Hermite jets, resultants and discriminants,
ordered root geometry (root bounds, Gauss--Lucas, Jensen disks, a sharp
Schoenberg second-moment inequality with its equality case), real-closed-field
transfer with Hermite signature counting, exact root loci under coefficient
uncertainty in both geometries, Newton profiles over an arbitrary ordered
exponent group, image balls and branch values, exact Rolle counts and the
finite root-cluster tree, support-controlled coprime Hensel factorization with
a fixed-domain parameter form that shrinks no neighbourhood, and universal
residue duality with the trace identity.  \emph{Three perturbation theorems
with three different hypotheses are kept separate rather than collapsed, as
are two discriminant layer formulas and two mapping theorems, with their
agreement proved; the alternative was to lose what each source reached that
the others did not.}}

\entry{Trigonometry on the surcomplex plane}
{surcomplex/trigonometry}
{article.pdf}{surcomplex\_trigonometry.zip and two same-named archives}
{Its organizing theorem is that every surcomplex direction has a \emph{finite}
surreal angle: $\mathrm{cis}:(\mathcal{O}_\R,+)\to(T(\No),\cdot)$ is surjective
with kernel exactly $2\pi\mathbb{Z}$, so
$\mathcal{O}_\R/2\pi\mathbb{Z}\cong T(\No)\cong T(\R)\times(\mathfrak{m}_\R,+)$,
and the whole of plane, spherical and hyperbolic triangle geometry is available
at arbitrary surreal length scale without ever choosing a value for
$\sin\omega$.  From it: canonical sine and cosine by Hahn summation, the unit
circle with its branch conventions and projective half-angle chart, inverse
functions on their full surreal domains, the angle-sum, cosine, sine, tangent,
half-angle, Heron, incircle, Euler, angle-bisector, cevian, Ceva,
inscribed-angle and Ptolemy theorems, an exact valuation dictionary between
sides, angles and area with a quadratic triangle-inequality defect formula and
two flat-triangle theorems, a rank-two collision-stable intersection algebra,
the canonical complex extension to the strip where sine is surjective with
completely determined fibres, finite Fourier inversion, Parseval and a
surcomplex Fej\'er--Riesz factorization, and the classification of all global
phase extensions, with the Ehrlich--Kaplan normalization placed inside that
classification.  Two coherence obstructions close it: there is no coherent
bounded sine of infinite frequency, and all-scale coherence forces a
polynomial.  \emph{No mathematical disagreement was found between the three
sources; the duplication was near-total on the first two thirds and the
divergence near-total on the last.}}



\entry{Differential algebra and differential equations over the surreal and surcomplex numbers}
{surcomplex/differential-equations}
{article.pdf}{surreal\_differential\_algebra.zip and six companion archives}
{Seven manuscripts on the Berarducci--Mantova derivation $\der$ on $\No$,
normalized by $\der\omega=1$, and its unique extension to $\SC$ --- unique
because differentiating $i^2+1=0$ forces $\der i=0$.  The organizing identity is
exact: the image of the logarithmic derivative is $\No+i\mathcal{A}$ with
$\mathcal{A}=\der\mathfrak{m}=\der\mathcal{O}$, so $\der y=(a+ib)y$ has a nonzero
surcomplex solution exactly when $b$ has a \emph{finite} primitive.  Two things
carry the subject.  The two printed forms of that image agree only through the
proved lemma $\der\mathcal{O}=\der\mathfrak{m}$, never as notational variants.
And $\mathcal{A}$ is a \emph{strict} subclass of $\mathfrak{m}$: $\omega^{-1}$ is
infinitesimal while its primitive $\log\omega$ is infinite, so
$\der y=i\omega^{-1}y$ has only the zero solution --- collapsing
$\der\mathcal{O}$ to $\mathcal{O}$ would make it solvable, which all seven
sources refute.  Hence there is no formal oscillator: $\der y=iy$ needs a finite
primitive of $1$, and those are $\omega+c$, all infinite, so $\der^2y+y=0$
factors as $(\der-i)(\der+i)y=0$ and has only $y=0$.  Also: the purely infinite
part of the accumulated phase as a complete gauge invariant in a four-term exact
sequence; a sharp threshold over real-exponent Hahn fields, solvable iff
$\min\operatorname{supp}b>1$, with the strictness at exponent~$1$ forced by the
$\log\omega$ term and a counterexample pair showing the criterion cannot be
globalized by valuation; Picard--Vessiot extensions with diagonal phase tori; and
a coherent nonlinear initial-value theorem.  \emph{Merged; the hazard was a
word.}  Three different objects were all called ``phase'' --- the finite angle in
$\mathcal{O}$, the infinitesimal residue in $\mathfrak{m}$ left after removing a
constant unit, and the accumulated phase, which is no angle at all and leaves
$\mathcal{O}$ whenever its purely infinite part is nonzero.  One source printed
two of them twenty lines apart.  The merged report fixes three symbols, relates
them rather than identifying them, and prints the tempting false sentence beside
the true one.  The letter $J$, which meant three unrelated things across the
batch, is retired.}

\entry{Rank-one non-Archimedean analytic geometry inside the surcomplex numbers}
{surcomplex/rank-one-berkovich}
{article.pdf}{surcomplex\_rank\_one\_berkovich.zip}
{The one report in the collection where convergence is \emph{real} convergence.
Inside the rank-one workspace $K=\C((t^\R))$ the absolute value
$|a|_v=\exp(-v(a))$ is genuine and $K$ is spherically complete, complete and
algebraically closed, so it is a legitimate base for Tate algebras and Berkovich
spaces.  Proves strict inclusions among the Tate, polynomial-coefficient Hahn,
common-domain, radius-free germ, formal-coefficient and open-disk analytic rings
with named witnesses, extending the ring table of
\nolinkurl{surcomplex/analytic-geometry}; gives a constructive one-variable
preparation theorem, principal-ideal and finite-quotient results, analytic root
counts and a valuation Rouch\'e theorem; constructs the Berkovich disk points and
the annulus retraction; and contrasts the contractible annulus spectrum with its
nonzero analytic de Rham class.  The worked case study is the entire series
$F(Z)=\sum_{n\ge0}t^{n^2}Z^n$, taken from \nolinkurl{foundations-and-computation/foundations}
and credited there: its full zero geometry is settled, and the first correction of
the $m$th normalized root occurs at $t^{m(m+1)}$.  \emph{This does not contradict
the no-convergence theorems of \nolinkurl{surcomplex/analysis}.}  The convergence
here is in the intrinsic valuation topology of one fixed rank-one field, not the
fine topology induced from $\SC$; and at $Z=t^{-\omega}$ the exponents $n^2-n\omega$
strictly decrease, so the support is not well ordered and $F$ is not an
all-surcomplex entire function --- the obstruction agrees with all-scale rigidity
rather than refuting it.  Written against this repository; its claim that a
Berkovich \emph{comparison} already exists in \nolinkurl{surcomplex/contours-and-stokes}
is correct, and what was absent is the development.}

\entry{Surcomplex spectral theory: singular values, infinitesimal rank, multiscale stability}
{surcomplex/spectral-theory}
{article.pdf}{surcomplex\_spectral\_theory.zip}
{The finite-matrix chapter the collection lacked, and the one report here with no
derivation, no phase and no differential equation in it.  Proves the finite
spectral theorem and the singular value decomposition over a real closed field
and its complexification \emph{without compactness}: real closedness suffices,
and the proofs exhibit the extremal witnesses rather than extracting them from a
compactness argument, so Courant--Fischer and Ky Fan hold with extrema
\emph{attained}.  On that base: positive square roots and polar decomposition,
generalized Hermitian eigenproblems, Moore--Penrose inverses, least squares, best
low-rank approximation, and conditioning.  The scale-sensitive core identifies
singular-value valuations with successive differences of minimum minor
valuations, and recovers both those valuations and the leading singular
coefficients from positive Gram coefficients one scale block at a time.
Applications include an intrinsic rank filtration, root-free scale elimination,
perturbation precision certificates, nonnormal pseudospectra, gap-sensitive
spectral projectors and Schur-complement corrections.  Written against this
repository; its gap claim checks out --- \nolinkurl{surcomplex/polynomial-algebra}
does carry Schur-type arguments and Hermite signatures, and no report in the
collection had a singular value, a pseudoinverse or a matrix polar decomposition
anywhere.}


\section{Foundations and computation}

Two reports that take the rest of the collection as their object.  Neither
belongs under \dir{surreal} or \dir{surcomplex}, because both range over the
whole subject.  Each is a merge of three manuscripts delivered together on
21~September 2026, and in each the hazard was not overlap but a collision of
near-identical names.

\entry{Foundations for surreal and surcomplex mathematics}
{foundations-and-computation/foundations}
{article.pdf}{surreal\_surcomplex\_foundations.zip + surcomplex\_foundations.zip
+ surreal\_foundations.zip}
{In what foundation does this mathematics legally live, and how would a proof
assistant encode it?  Compares ZF/ZFC with definable classes, GB/GBC/NBG,
Kelley--Morse, Grothendieck universes and TG, Feferman reflection, constructive
set theory, dependent type theory, univalent foundations and native Lean, and
proves the size obstructions each must respect: no unrestricted cut operation
can accept every subclass of its own carrier, the surreals are not Dedekind
complete, and a universe cannot cut itself.  The positive counterpart is a
localization theorem --- every \emph{set} of surcomplex numbers already lies in
one divisible set-sized workspace $K_\Gamma=\C((t^\Gamma))$ --- with a
finite-algebra descent theorem explaining why a finite zero scheme gains no
points in a larger workspace.  Carries two series over $\C((t^\Q))$ that differ
only in the sign of an exponent and reach opposite conclusions:
$\sum t^{-n^2}X^n$ is summable at no nonzero argument, while $\sum t^{+n^2}z^n$
is summable everywhere and is a \emph{nonpolynomial} function coherent at every
radius --- which does not refute the all-scale rigidity theorem of
\nolinkurl{surcomplex/analysis}, stated for class functions on the whole $\SC$,
but shows it does not transfer to a fixed value group.  \emph{Merged; the danger
was aggregation.  All three manuscripts read the same Lean repository at the
same revision on the same day, so their agreement is one act of source reading
performed three times, not three confirmations.  The findings are printed once
and the verification boundary has its own early section.}  Three of the five
shipped Lean files were compiled for this repository and \texttt{\#print axioms}
run on every declaration: no \texttt{sorryAx} and no \texttt{Classical.choice}
anywhere, and both size-obstruction theorems depend on no axioms at all.  The
toolchain used was not the one the sources pin, and one shipped file does not
compile as delivered.}

\entry{Surreal and surcomplex numbers in symbolic computer algebra}
{foundations-and-computation/computer-algebra}
{article.pdf}{surreal\_surcomplex\_cas.zip and two same-named archives}
{What can actually be computed exactly, and what only approximated or denoted?
Separates finite exact denotation, effective coefficient access, decidable
equality and order, and certified approximation into capability tiers, then
works through an ordered rational-monomial core, real-algebraic closure,
certified Hahn series, formal germs, analytic lifts, transseries, Newton
profiles, Hensel lifting, finite quotient algebras, residues, and a proposed
Wolfram Language architecture.  \emph{Merged; the danger was three prototypes
with one name.}  Each manuscript ships an executable Wolfram package, two of
them near-anagrams of each other, two exporting six identical public symbol
names with incompatible signatures, and all three calling themselves the
rational monomial core.  They are three different objects: coefficients $\Q$ in
one and $\Q(i)$ in the others; exponent lattices $\Z^d$ against $\Q^r$, so a
ramified root is representable in one and not the next; and two of them order
their monomial variables by opposite conventions, which reverses every
inequality.  The report never writes ``the prototype'' and tabulates the
divergences instead.  All three check suites were rerun for this repository and
reproduce exactly --- 34/34, 110/110 and 88/88 in Wolfram, 69/69 and 323/323 in
Python --- but one package cannot be loaded at all as shipped: an unclosed
bracket stops it parsing, which is why its own record submits definitions rather
than loading the file.}


\end{document}
